%


\section{Solvability and right groups}
\label{sec:solv}
Throughout this work, we consider Lie algebras over $\C$. In general, if $\?g?$ is a Lie 
algebra, its {\sl derived algebra,} $\?g?'=[\?g?,\?g?]$, 
is the linear span of all products $[x,y]$, $x,y\in\?g?$. 
This is an ideal of $\?g?$. Continuing, one defines a 
sequence of ideals 
$$\?g?\supset\?g?' \supset \?g?'' \supset \cdots \supset \?g?^{(n)} \supset \cdots$$ with
$\?g?^{(n+1)}=[\?g?^{(n)},\?g?^{(n)}]$,
which is the {\sl derived series\/} of $\?g?$. 
If this series terminates at $0$ after finitely many steps, then $\?g?$ is {\sl solvable}. 
Note that the last nonzero term of the series  of a solvable algebra is abelian.\\

The main property of a solvable Lie algebra
that we are using here is {\em Lie's Theorem\/} to the effect that (over an algebraically closed field)
a solvable Lie algebra of matrices can be simultaneously upper-triangularized 
(\cite[p.\tsp 50]{JA}, also see \cite[p.\tsp 21--26]{RR}).
In this context the work by Radjavi \& Rosenthal \cite{RR} is of particular interest. Especially,
Radjavi's Theorem on permutable traces \cite[p.\tsp 33]{RR} 
is part of the inspiration behind the proof.\\

There are three Lie algebras associated with a coloring $\Co$ of 
a strongly connected digraph with coloring semigroup $S$ and its kernel $K$.
In general, we denote by $\?g?(\cdot)$ the Lie algebra generated by 
transformations from a given set. In this context, we specify
$\cS=\?g?(S)$ the Lie algebra generated by the full semigroup. The underlying set is the 
same as the semigroup algebra (without 1) with commutator as product. The Lie algebra
$\L=\?g?(\Co)$, generated by the transformations in $\Co$, is a subalgebra of $\cS$ and is our main
object of interest. Finally, $\cK=\?g?(K)$, is a subalgebra of $\cS$, that is, in fact, a Lie ideal of $\cS$
since $K$ is an ideal in the semigroup. 
The main feature is that all three Lie algebras are solvable
if and only if the generators $\Co$ are simultaneously triangularizable if and only if $\L$ is
solvable. In that case, we call the graph {\it Lie-solvable\/}.\\

\begin{theorem} \label{thm:RG} \label{thm:4.14} 
If $\L$ is solvable, then the kernel $K$ is a right group.
\end{theorem}

We start with a Lemma.
\begin{lemma}\label{lem:dirprod} \label{lem:4.1}    
If $\L$ is solvable, then the kernel $K$ is isomorphic 
to a Rees product semigroup that is a direct product.
\end{lemma}

\begin{proof}
We want to show that
$$
  K \cong X \times G \times Y \hbox{ \ where if \ }
	(x_1,g_1,y_1), (x_2,g_2,y_2) \in X \times G \times Y
$$
then
$$
  (x_1,g_1,y_1) (x_2,g_2,y_2) = (x_1x_2,g_1g_2,y_1y_2).
$$
It is well known that this is equivalent to the property: \\
\begin{center}
 for any two idempotents $e,f \in K$, $(ef)^2=ef$
\end{center}
\hfill\break
In other words, the idempotents of $K$ should form a {\sl band,} a semigroup of idempotents.\\

Note that if we show that $\rk(ef) =  \tr(ef)$, then we are done, since for a function 
$$
  \rk(f) = \tr(f) \hbox{ iff } 
	f \hbox{ is an idempotent.}
$$
Let $e,f\in K$, then,
since $ef$ is contained in some maximal subgroup of $K$,
$(ef)^p$ is an idempotent for some $p\ge1$.
Since $\L$ is solvable, the matrices of $\L$ can be made simultaneously upper triangular, thus $\tr((ef)^p) = \tr(e^pf^p)$.
Also, since both $(ef)^p$ and $ef$ are in $K$, $\rk((ef)^p) = \rk(ef)$. 
Thus we have $\rk(ef) = \rk((ef)^p) = \tr((ef)^p) = \tr(e^pf^p) = \tr(ef)$, since $e^p=e$ and
$f^p=f$.\\
Therefore, $ef$ is an idempotent and it follows that the Rees product is a direct product.
\qed
\end{proof}

Now we proceed with the proof of the theorem.\\

\begin{proof}
Consider two partitions $\pi_1,\pi_2$ of the vertices $V$, corresponding to minimal right ideals of $K$.
Let $\pi_1 = \{P_1,P_2,\ldots,P_r\}$, $\pi_2 = \{Q_1,Q_2,\ldots,Q_r\}$, and suppose $R_1,R_2,\ldots,R_t$ are all possible ranges of the functions in $K$. (Note: $r$ is the rank of the kernel.)\\

Label the idempotents $e_i = E\big(H(\pi_1,R_i)\big)$ and $f_i = E\big(H(\pi_2,R_i)\big)$ where $H(\pi_j,R_i)$ is the $H$-class determined by partition $\pi_j$ and range $R_i$.
Notice that since $K$ is a direct product, $e_if_j = e_j$, since $e_if_j \in H(\pi_1,R_j)$ and is an idempotent.\\

Let $R_1 = \{x_1,x_2,\ldots,x_r\}$, with $x_i \in P_i$.
Consider the sets  $O_j = \{x_je_i \mid i=1,\ldots,t\}$.
Each idempotent acts as a set identity function on its range, so that 
$x_je_i \in P_j$ for $i = 1,\ldots,t$ and $j=1,\ldots,r$.
Hence $O_j \subset P_j$ for $j = 1,\ldots,r$.
In particular, $O_{j_1} \cap O_{j_2} = \emptyset$, if $j_1 \not= j_2$.\\

Notice that if, for some $x\in P_j$, $xe_i=x$, then $x_je_i=x$ will pick up $x\in O_j$.
Thus,
$$
  V = \bigcup^r_{j=1} O_j
$$
since $\bigcup\limits^t_{i=1} R_i = V$ so that
every element of $V$ is in the range of some idempotent $e_i$.
So, in fact, $O_j = P_j$ for $j = 1,\ldots,r$.\\

Similarly, $Q_j = \{ x_jf_k \mid k = 1,\ldots,t\}$.
But for each $k$, $1\le k\le t$,
$$
  x_jf_k = x_je_1f_k = x_je_k \in P_j.
$$
Thus $Q_j \subset P_j$ for $j = 1,\ldots,r$.
Then, since $V=\bigcup\limits^r_{j=1} Q_j = \bigcup\limits^r_{i=1} P_i $, it is clear that 
 $Q_j = P_j$ for $j = 1,\ldots,r$ and thus, $\pi_1 = \pi_2$.
Hence $K$ is a right group and the theorem is proved.
\qed
\end{proof}

Combining this with Theorem \ref{thm:4.13}, we have 

\begin{corollary}\label{cor:sync}
If the Lie algebra $\L$ of a coloring is solvable, then there exists a coloring semigroup
of the graph that contains a synchronizing instruction. In other words, 
the road coloring conjecture holds for Lie-solvable graphs.

\end{corollary}

